% Рубежный контроль №1 «Моделирование» — Вариант 2 (по фото пользователя)
% Вопросы 1–6: со страницы с заголовком «Вариант 2»
% Вопросы 7–10: со страницы с СМО/методом моментов/соцопросом
\documentclass[12pt,a4paper]{article}
\usepackage[T2A]{fontenc}
\usepackage[utf8]{inputenc}
\usepackage[russian]{babel}
\usepackage{amsmath,amssymb}
\usepackage{geometry}
\geometry{margin=2.2cm}
\usepackage{enumitem}
\setlist{nosep,leftmargin=*}

\title{Рубежный контроль №1 «Моделирование»\\Вариант 2 — ответы}
\author{}
\date{}

\begin{document}
\maketitle

\begin{enumerate}[label=\textbf{\arabic*.}]

\item \textbf{Имитационная модель (определение).}

Имитационная модель — численный метод вычислительного эксперимента, позволяющий имитировать поведение реальной системы на ЦВМ в заданный или формируемый период времени при заданных входных данных (и управляющих воздействиях). Для многокомпонентных систем используют модельное время и схемы квазипараллелизма/синхронизации.

\item \textbf{Феноменологическая модель (определение).}

Феноменологическая модель — модель, построенная на результатах конкретного (уже завершённого) эксперимента и описывающая наблюдаемую закономерность без явного вывода из физического механизма; область применимости ограничена условиями эксперимента.

\item \textbf{Натурный эксперимент (определение).}

Натурный эксперимент — эксперимент с реальными объектами и процессами (в естественных/реальных условиях функционирования), в отличие от вычислительного эксперимента на математической модели.

\item \textbf{Классификация сложных динамических систем (перечислить).}

В курсе выделяют (в т.\,ч.) следующие классы динамических систем:
\begin{enumerate}[label=\arabic*)]
\item простые динамические системы;
\item структурно-сложные динамические системы;
\item сложные динамические системы, меняющие поведение во времени;
\item структурно-сложные гибридные системы.
\end{enumerate}

\item \textbf{Выбор модельного времени с переменным шагом.}

Модельное время — дискретное время в единицах реального времени, обеспечивающее внутреннюю и внешнюю синхронизацию компонент модели.
При \textbf{переменном шаге} кванты модельного времени адаптируются к процессу:
\begin{itemize}
\item если два и более событий попадают в один и тот же момент, выполняют последовательное описание компонент и корректируют шаг;
\item если длительное время нет событий, шаг увеличивают (чтобы избежать «простоя» моделирования).
\end{itemize}

\item \textbf{Методика построения аналитической динамической модели.}

Методика (кратко):
\begin{enumerate}[label=\arabic*)]
\item по постановке задачи строят схему процесса во времени ($t_0$, $t^{*}$, $T$);
\item уточняют входные/выходные данные и промежуточные переменные;
\item выделяют физико-химические воздействия, влияющие на выход;
\item упрощают модель (отбрасывая малозначимые воздействия);
\item формализуют модель системой уравнений/неравенств, связывающих входы и выходы;
\item детализируют каждое воздействие и уточняют коэффициенты/константы;
\item проверяют согласованность: число неизвестных не превышает числа уравнений, заданы начальные/краевые условия.
\end{enumerate}
\item \textbf{Пример аналитической модели системы массового обслуживания.}

Обозначение \textbf{$M/M/1$} читается так:
\begin{itemize}
\item первая $M$ — \emph{Markovian} входной поток: интервалы между заявками экспоненциальны, поток пуассоновский с интенсивностью $\lambda$;
\item вторая $M$ — \emph{Markovian} обслуживание: времена обслуживания экспоненциальны с интенсивностью $\mu$;
\item $1$ — число каналов обслуживания (серверов).
\end{itemize}
Аналогично \textbf{$M/M/c$} — то же самое, но каналов $c$.

\textbf{1) Модель $M/M/1$ (пример с вычислениями).}
\[
\lambda=2\ \text{заявки/мин},\qquad \mu=3\ \text{заявки/мин}.
\]
Коэффициент загрузки:
\[
\rho=\frac{\lambda}{\mu}=\frac{2}{3}<1 \quad (\text{устойчивость}).
\]
Стационарное распределение числа заявок в системе:
\[
P\{N=n\}=(1-\rho)\rho^n,\quad n=0,1,2,\dots,
\]
Среднее число заявок в системе:
\[
L=\mathbb{E}N=\frac{\rho}{1-\rho}=\frac{\tfrac23}{1-\tfrac23}=\frac{2/3}{1/3}=2.
\]
Среднее время пребывания заявки в системе:
\[
W=\mathbb{E}T=\frac{1}{\mu-\lambda}=\frac{1}{3-2}=1\ \text{минута}.
\]
Среднее время ожидания в очереди:
\[
W_q=W-\frac{1}{\mu}=1-\frac{1}{3}=\frac{2}{3}\ \text{мин} \approx 0.667\ \text{мин}.
\]

\textbf{2) Модель $M/M/c$ (возьмём $c=2$, пример с вычислениями).}
\[
\lambda=3\ \text{заявки/мин},\qquad \mu=2\ \text{заявки/мин},\qquad c=2.
\]
Загрузка системы:
\[
\rho=\frac{\lambda}{c\mu}=\frac{3}{2\cdot 2}=0.75<1.
\]
Обозначим $\alpha=\lambda/\mu=1.5$. Тогда
\[
P_0=\left(\sum_{n=0}^{c-1}\frac{\alpha^n}{n!}+\frac{\alpha^c}{c!}\cdot \frac{1}{1-\rho}\right)^{-1}
=\left(1+\alpha+\frac{\alpha^2}{2}\cdot\frac{1}{1-\rho}\right)^{-1}.
\]
Подстановка чисел:
\[
P_0=\left(1+1.5+\frac{2.25}{2}\cdot\frac{1}{0.25}\right)^{-1}
=\left(2.5+1.125\cdot 4\right)^{-1}
=\frac{1}{7}.
\]
Вероятность ожидания (формула Эрланга~C):
\[
P_{\text{wait}}=\frac{\alpha^c}{c!}\cdot\frac{1}{1-\rho}\cdot P_0
=\frac{2.25}{2}\cdot 4\cdot \frac{1}{7}
=\frac{4.5}{7}\approx 0.643.
\]
Средняя длина очереди:
\[
L_q=P_{\text{wait}}\cdot \frac{\rho}{1-\rho}
=\frac{4.5}{7}\cdot \frac{0.75}{0.25}
=\frac{4.5}{7}\cdot 3
=\frac{13.5}{7}\approx 1.929.
\]
Среднее время ожидания:
\[
W_q=\frac{L_q}{\lambda}\approx \frac{1.929}{3}\approx 0.643\ \text{мин},
\]
а среднее время в системе:
\[
W=W_q+\frac{1}{\mu}\approx 0.643+0.5=1.143\ \text{мин}.
\]

\item \textbf{Построение модифицированного метода моментов (линейная функциональная зависимость).}

\textbf{Постановка.} Даны пары наблюдений $(x_i,y_i)$, $i=1,\dots,n$. Предполагается линейная зависимость среднего отклика:
\[
\mathbb{E}[Y\mid X=x]=\theta_0+\theta_1 x.
\]
Модифицированный метод моментов для оценки параметров связи строит систему уравнений, приравнивая теоретические характеристики выборочным аналогам:
\[
\mathbb{E}[Y]=\theta_0+\theta_1\mathbb{E}[X],\qquad
\mathrm{cov}(X,Y)=\theta_1\,\mathrm{var}(X).
\]
Заменяя теоретические величины выборочными:
\[
\bar x=\frac1n\sum_{i=1}^n x_i,\quad
\bar y=\frac1n\sum_{i=1}^n y_i,\quad
\widehat S_{xx}=\frac1n\sum_{i=1}^n (x_i-\bar x)^2,\quad
\widehat S_{xy}=\frac1n\sum_{i=1}^n (x_i-\bar x)(y_i-\bar y),
\]
получаем оценки параметров:
\[
\hat\theta_1=\frac{\widehat S_{xy}}{\widehat S_{xx}},\qquad
\hat\theta_0=\bar y-\hat\theta_1\bar x,
\]
при $\widehat S_{xx}>0$.

\item \textbf{9–10. Имеются результаты соцопроса 10 человек. Оценить силу связи между показателями.}

Пусть ответы «Да/Нет» кодируются как $1/0$. Для связи двух дихотомических признаков будем использовать \textbf{коэффициент $\varphi$} — это \textbf{коэффициент Пирсона} для таблицы сопряжённости $2\times2$ (он же $\varphi=\sqrt{\chi^2/n}$ с учётом знака). \textbf{Коэффициент Юла $Q$} — альтернативная мера ассоциации для $2\times2$; здесь используем именно $\varphi$ как наиболее стандартный «корреляционный» показатель.
\[
\varphi=\frac{ad-bc}{\sqrt{(a+b)(c+d)(a+c)(b+d)}},
\]
где $a$ — число пар $(1,1)$, $b$ — $(1,0)$, $c$ — $(0,1)$, $d$ — $(0,0)$.

\textbf{Связь между вопросами 1 и 2.}

\textbf{Шаг 1. Кодирование.} «Да» $\to 1$, «Нет» $\to 0$. Обозначим через $Q_1,Q_2,Q_3,Q_4$ \emph{вопросы анкеты} (это не коэффициент Юла; это просто имена переменных «ответ на вопрос №1», «ответ на вопрос №2» и т.\,д.). По фото (10 человек):
\[
Q_1=(1,0,1,0,0,0,0,1,0,1),\quad
Q_2=(1,1,0,1,1,0,1,0,0,0).
\]

\textbf{Шаг 2. Таблица сопряжённости $2\times2$.}
Подсчёт пар $(Q_1,Q_2)$ даёт:
\[
a=\#(1,1)=1,\quad b=\#(1,0)=3,\quad c=\#(0,1)=4,\quad d=\#(0,0)=2.
\]
То есть
\[
\begin{array}{c|cc}
 & Q_2=1 & Q_2=0\\ \hline
Q_1=1 & a=1 & b=3\\
Q_1=0 & c=4 & d=2
\end{array}
\]

\textbf{Шаг 3. Подстановка в формулу.}
\[
\varphi_{12}=\frac{ad-bc}{\sqrt{(a+b)(c+d)(a+c)(b+d)}}
=\frac{1\cdot2-3\cdot4}{\sqrt{(1+3)(4+2)(1+4)(3+2)}}
=\frac{-10}{\sqrt{600}}
\approx -0.408.
\]

\textbf{Ответ:} $\varphi_{12}\approx -0.408$.
\[
\varphi_{12}\approx -0.408.
\]

\textbf{Связь между вопросами 3 и 4.}

\textbf{Шаг 1. Кодирование.} По фото:
\[
Q_3=(0,0,0,1,0,0,0,0,1,0),\quad
Q_4=(0,1,0,0,1,1,0,0,0,1).
\]

\textbf{Шаг 2. Таблица сопряжённости.}
\[
a=\#(1,1)=0,\quad b=\#(1,0)=2,\quad c=\#(0,1)=4,\quad d=\#(0,0)=4,
\]
то есть
\[
\begin{array}{c|cc}
 & Q_4=1 & Q_4=0\\ \hline
Q_3=1 & a=0 & b=2\\
Q_3=0 & c=4 & d=4
\end{array}
\]

\textbf{Шаг 3. Подстановка.}
\[
\varphi_{34}=\frac{0\cdot4-2\cdot4}{\sqrt{(0+2)(4+4)(0+4)(2+4)}}
=\frac{-8}{\sqrt{384}}
\approx -0.408.
\]

\textbf{Ответ:} $\varphi_{34}\approx -0.408$.
\[
\varphi_{34}\approx -0.408.
\]

\textbf{Остальные пары (аналогично).} По тем же правилам подсчёта $a,b,c,d$:
\begin{itemize}
\item $Q_1$ и $Q_3$: $(a,b,c,d)=(0,4,2,4)$,
\[
\varphi_{13}=\frac{0\cdot4-4\cdot2}{\sqrt{(0+4)(2+4)(0+2)(4+4)}}=\frac{-8}{\sqrt{384}}\approx -0.408.
\]
\item $Q_1$ и $Q_4$: $(a,b,c,d)=(1,3,3,3)$,
\[
\varphi_{14}=\frac{1\cdot3-3\cdot3}{\sqrt{(1+3)(3+3)(1+3)(3+3)}}=\frac{-6}{24}=-0.25.
\]
\item $Q_2$ и $Q_3$: $(a,b,c,d)=(1,4,1,4)$,
\[
\varphi_{23}=\frac{1\cdot4-4\cdot1}{\sqrt{(1+4)(1+4)(1+1)(4+4)}}=0.
\]
\item $Q_2$ и $Q_4$: $(a,b,c,d)=(2,3,2,3)$,
\[
\varphi_{24}=\frac{2\cdot3-3\cdot2}{\sqrt{(2+3)(2+3)(2+2)(3+3)}}=0.
\]
\end{itemize}

\textbf{Итог по всем парам.} (знак «$-$» — отрицательная связь, «$0$» — связи не выявлено):
\[
\varphi_{12}\approx -0.408,\quad
\varphi_{13}\approx -0.408,\quad
\varphi_{14}=-0.25,\quad
\varphi_{23}=0,\quad
\varphi_{24}=0,\quad
\varphi_{34}\approx -0.408.
\]
Вывод: наиболее заметная связь (умеренная отрицательная) наблюдается для пар $12,13,34$; для $14$ связь слабая отрицательная; для $23$ и $24$ линейной ассоциации по $\varphi$ не выявлено.

\end{enumerate}
\end{document}
